{
\renewcommand{\arraystretch}{1.5}
\setlist[enumerate]{leftmargin=*, labelsep=0.5em}

\begin{tabularx}{\textwidth}{|>{\raggedright\arraybackslash}X|>{\raggedright\arraybackslash}X|}
\hline
\textbf{Tên Use case} & Chỉnh sửa bộ thẻ ghi nhớ \\
\hline
\textbf{ID} & UC-0033 \\
\hline
\textbf{Tác nhân} & Người dùng đã xác thực (chủ sở hữu bộ thẻ) \\
\hline
\textbf{Mô tả} & Người dùng chỉnh sửa thông tin bộ thẻ ghi nhớ \\
\hline
\textbf{Điều kiện tiên quyết} & Bộ thẻ tồn tại, người dùng là tác giả \\
\hline
\textbf{Điều kiện sau} & Thông tin bộ thẻ được cập nhật thành công \\
\hline
\textbf{Kích hoạt} & Người dùng chọn chỉnh sửa bộ thẻ \\
\hline
\textbf{Luồng chính} &
\begin{enumerate}
    \item Người dùng mở bộ thẻ ghi nhớ.
    \item Người dùng nhấn Chỉnh sửa.
    \item Người dùng chỉnh sửa: tên, mô tả, chế độ hiển thị (công khai/riêng tư), thẻ phân loại.
    \item Người dùng nhấn Lưu.
    \item Hệ thống cập nhật bộ thẻ và xác nhận thành công.
\end{enumerate} \\
\hline
\textbf{Luồng thay thế} &
\begin{enumerate}
    \item Không có luồng thay thế.
\end{enumerate} \\
\hline
\textbf{Luồng ngoại lệ} &
\begin{enumerate}
    \item \textbf{Bộ thẻ không tồn tại:}
    \begin{enumerate}
        \item[1Err.] Hệ thống trả về lỗi 404.
    \end{enumerate}
    \item \textbf{Không có quyền:}
    \begin{enumerate}
        \item[1Err2.] Hệ thống trả về lỗi 403 khi người dùng không phải tác giả.
    \end{enumerate}
\end{enumerate} \\
\hline
\end{tabularx}
\captionof{table}{Đặc tả Use case: UC-0033 Chỉnh sửa bộ thẻ ghi nhớ}
}
